\subsection{Perancangan Sistem}

Tahap perancangan sistem bertujuan menggambarkan alur kerja sistem \textit{checkout} otomatis Nasi Padang mulai dari pelatihan model hingga penggunaan model pada proses \textit{checkout}. Perancangan dibagi menjadi dua bagian, yaitu rancangan pelatihan model YOLO11-seg dan rancangan proses \textit{checkout}.

Rancangan pelatihan model menjelaskan alur penggunaan dataset yang telah dianotasi dan dibagi menjadi data pelatihan, validasi, dan pengujian. Data pelatihan digunakan untuk melatih YOLO11-seg melalui proses prediksi, perhitungan \textit{loss}, dan pembaruan bobot model. Performa model dipantau menggunakan data validasi hingga diperoleh model akhir yang digunakan dalam proses pengujian dan inferensi.

Rancangan proses \textit{checkout} menjelaskan alur pengolahan citra piring yang diperoleh dari kamera. Model YOLO11-seg menghasilkan prediksi kelas dan \textit{segmentation mask} untuk setiap instans objek makanan. Setiap instans kemudian dihitung berdasarkan kelasnya dan dicocokkan dengan tabel harga tetap untuk memperoleh subtotal setiap kelas serta total harga keseluruhan.

Rancangan pelatihan model dan proses \textit{checkout} dijelaskan lebih lanjut pada subbagian berikut.

\subsubsection{Rancangan Pelatihan Model}

Pelatihan dilakukan menggunakan YOLO11-seg dengan bobot pralatih sebagai titik awal. Bobot tersebut selanjutnya disesuaikan terhadap sembilan kelas makanan pada dataset penelitian sehingga proses pelatihan tidak dimulai dari bobot acak.

Data pelatihan diberikan kepada model dalam bentuk citra beserta \textit{ground truth} yang memuat label kelas dan poligon segmentasi setiap instans. Pada saat data dimuat, citra disesuaikan dengan ukuran masukan model dan dapat menerima augmentasi sesuai konfigurasi pelatihan.

Selama pelatihan, model menghasilkan prediksi yang dibandingkan dengan \textit{ground truth} untuk memperoleh nilai \textit{loss}. Nilai tersebut digunakan dalam proses \textit{backpropagation} dan pembaruan bobot. Proses ini dilakukan berulang selama sejumlah \textit{epoch} yang ditetapkan pada tahap implementasi.

Data validasi digunakan selama pelatihan untuk memantau perkembangan performa model tanpa digunakan untuk memperbarui bobot. Hasil validasi digunakan sebagai dasar pemilihan bobot model yang selanjutnya digunakan pada evaluasi akhir.

Model terpilih kemudian dievaluasi menggunakan data pengujian yang tidak digunakan untuk pembaruan maupun pemilihan bobot. Evaluasi dilakukan terhadap kemampuan model dalam mendeteksi, mengklasifikasikan, membentuk \textit{segmentation mask}, serta menghasilkan jumlah instans yang sesuai dengan \textit{ground truth}. Varian YOLO11-seg, ukuran masukan, jumlah \textit{epoch}, \textit{batch size}, dan konfigurasi pelatihan lainnya dicatat sesuai implementasi yang benar-benar digunakan.

\subsubsection{Rancangan Proses Checkout}

Proses \textit{checkout} dimulai ketika piring yang berisi makanan ditempatkan pada area pengambilan citra. Kamera yang berada pada posisi tetap di atas area tersebut memperoleh citra piring dalam kondisi pencahayaan yang relatif terkontrol. Citra kemudian diberikan sebagai masukan kepada model YOLO11-seg yang telah melalui proses pelatihan.

Pada tahap inferensi, YOLO11-seg memproses citra dan menghasilkan prediksi untuk setiap instans objek makanan. Setiap prediksi memuat lokasi objek, nilai kepercayaan, label kelas, dan \textit{segmentation mask}. Prediksi dengan nilai kepercayaan yang memenuhi batas yang ditentukan dipertahankan sebagai hasil pengenalan.

Label kelas digunakan untuk menentukan jenis objek makanan, sedangkan \textit{segmentation mask} digunakan untuk memisahkan setiap instans yang terdapat pada piring. Setiap hasil prediksi yang dipertahankan diperlakukan sebagai satu instans. Apabila terdapat lebih dari satu instans dari kelas yang sama, seluruh instans tersebut dihitung dan dikelompokkan berdasarkan kelasnya.

Hasil penghitungan kemudian diteruskan ke program kalkulasi harga. Program menggunakan tabel harga tetap yang memuat pasangan antara kelas makanan dan harga setiap instans. Subtotal setiap kelas dihitung dengan mengalikan jumlah instans yang dikenali dengan harga pada kelas tersebut. Seluruh subtotal selanjutnya dijumlahkan untuk memperoleh total harga.

Informasi yang dihasilkan pada proses \textit{checkout} terdiri atas kelas makanan yang dikenali, jumlah instans pada setiap kelas, harga setiap instans, subtotal setiap kelas, dan total harga keseluruhan. \textit{Segmentation mask} tidak digunakan untuk memperkirakan berat, volume, luas porsi, maupun kandungan gizi makanan.

Proses tersebut menghasilkan mekanisme \textit{checkout} otomatis yang menghubungkan hasil pengenalan YOLO11-seg dengan program kalkulasi harga. Sistem dirancang sebagai \textit{proof-of-concept} dan tidak mencakup proses pembayaran elektronik maupun pencetakan bukti pembayaran.

